\documentclass[a4paper,10pt]{article}
\usepackage[utf8]{inputenc}
\usepackage[T1]{fontenc}
\usepackage[french]{babel}
\usepackage{amsmath,amssymb}
\usepackage{geometry}
\usepackage{array,tabularx}
\usepackage{tikz}
\geometry{margin=1.6cm}
\pagestyle{empty}
\newcommand{\ligne}{\par\vspace{0.55cm}\noindent\dotfill\par}
\newcommand{\carreaux}[1]{\begin{center}\begin{tikzpicture}\draw[step=0.5cm,gray!35,very thin] (0,0) grid (17,#1);\end{tikzpicture}\end{center}}
\newenvironment{exercice}[2]{\paragraph{Exercice #1 \hfill #2 points}\mbox{}\par}{\medskip}
\begin{document}
\begin{center}\Large\bfseries Devoir surveillé - Statistiques et fonctions\end{center}
\vspace{2mm}
\begin{exercice}{1}{5}Une série a pour valeurs $4,7,8,8,10,13,16,18$. Calculer moyenne, médiane et étendue.\end{exercice}\begin{exercice}{2}{5}Pour la fonction $f(x)=2x-3$, calculer $f(0)$, $f(5)$ et l'antécédent de $7$.\end{exercice}\begin{exercice}{3}{5}Un graphique passe par les points $(0;1)$ et $(3;7)$. Déterminer si la situation peut être proportionnelle. Justifier.\end{exercice}\begin{exercice}{4}{5}Une enquête donne $40\%$ de réponses favorables sur $250$ personnes. Calculer l'effectif correspondant.\end{exercice}\end{document}
